%--------------------
% Packages
% -------------------
\documentclass[11pt,a4paper]{article}
\usepackage[utf8x]{inputenc}
\usepackage[T1]{fontenc}
%\usepackage{gentium}
\usepackage{mathptmx} % Use Times Font
\usepackage{soul}


\usepackage[pdftex]{graphicx} % Required for including pictures
\usepackage[english]{babel} % Swedish translations
\usepackage[xetex,linkcolor=black,pdfborder={0 0 0}]{hyperref} % Format links for pdf
\usepackage{calc} % To reset the counter in the document after title page
\usepackage{enumitem} % Includes lists

\frenchspacing % No double spacing between sentences
\linespread{1.2} % Set linespace
\usepackage[a4paper, lmargin=0.1\paperwidth, rmargin=0.1\paperwidth, tmargin=0.1\paperheight, bmargin=0.1\paperheight]{geometry} %margins
%\usepackage{parskip}

\usepackage[all]{nowidow} % Tries to remove widows
% \usepackage[protrusion=true,expansion=true]{microtype} % Improves typography, load after fontpackage is selected

\usepackage{lipsum} % Used for inserting dummy 'Lorem ipsum' text into the template

\usepackage{listings}
\usepackage{tabularx}

\lstdefinestyle{mystyle}{
    basicstyle=\ttfamily\footnotesize,
    breakatwhitespace=false,         
    breaklines=false,                                
    keepspaces=true,                 
    numbers=left,                    
    numbersep=5pt,                  
    showspaces=true,                
}

\lstset{style=mystyle}

%-----------------------
% Set pdf information and add title, fill in the fields
%-----------------------
\hypersetup{ 	
pdfsubject = {},
pdftitle = {sirius\_master\_thesis},
pdfauthor = {Eske Haack}
}

\setlength{\parskip}{10pt}
\setlength{\parindent}{0em}

\begin{document}

\section*{Master Thesis Project}

\textbf{English Title}:
Conditional Diffusion Models for Temperature and Precipitation Super-resolution over Europe

\textbf{Danish Title}:
Betingede diffusionsmodeller til superopløsning af temperatur og nedbør over Europa

\small
\textbf{Name}:
Eske Haack (s214643) \\
\textbf{Start date}:
2026-09-01 \\
\textbf{End date}:
2027-02-01 \\
\textbf{ECTS}:
30 Points

\normalsize
\subsection*{Description}
This master thesis will investigate the application of conditional Diffusion Models on spatial super-resolution. Specifically the focus will be on super-resolution of daily temperature and precipitation data over the European Domain, using low-resolution EC-Earth3-Veg climate simulations as predictors and high-resolution HCLIM simulations as targets. The downscaled fields will be evaluated for their ability to represent future climate indicators, including heave-wave-relevant temperature extremes and heavy-precipitation characteristics.

The research aims to investigate if diffusion models are a viable option for temperature and precipitation super-resolution, in a climate-projection setting. Diffusion models have been chosen for their generative properties, as well as the ability to compute ensembles, and thus uncertainties, of predictions.
In this study a series of diffusion-approaches and -architectures will be investigated to improve super-resolution performance. As well as a shorter investigation on the models transferability and across realizations. 

\subsection*{Timeline}
\begin{tabularx}{\textwidth}{lX}
    Week 1-4 & Literature study of relevant material. The goal is to create a general understanding of the climate modelling field and an overview of the current state-of-the-art machine learning approaches for climate super-resolution. \\
    Week 4 & \textbf{Milestone}: Hand-in of project plan. This should include a general description of the project, and an agreed list learning objectives. \\
    Week 5-13 & Developing the diffusion model. This time period will be used to build and study the modifications made to the diffusion model. \\
    Week 14-18 & Running experiments and evaluating results. \\
    Week 18-20 & Post-processing results, computing climatological features to evaluate the trained model against the physics-based regional climate model \\
    Week 20-22 & Finalize writing the report.
\end{tabularx}

\subsection*{Learning objectives}
\subsubsection*{Specific learning objectives for this thesis}
\begin{enumerate}
    \item Given a conditional diffusion model, the student can explain the function of the model and identify improvements. 
    \item Given a set of real-world data, the student can produce and analyse an appropriate machine learning model in a high-level programming language. 
    \item Given a machine learning model, the student can analyse and evaluate model performance by designing evaluation experiments.
\end{enumerate}

\subsubsection*{General learning objectives of the MSc in Engineering at DTU}
A graduate of the MSc programme from DTU:
\begin{enumerate}
    \item can identify and reflect on technical scientific issues and understand the interaction between the various components that make up an issue
    \item can, based on a clear academic profile, apply elements of current research at an international level to develop ideas and solve problems
    \item masters technical scientific methodologies, theories, and tools, and can take a holistic view of and delimit a complex, open issue, put it into a broader academic and societal perspective, and, on this basis, propose a variety of possible actions while considering sustainability
    \item can develop relevant models, systems, technologies, and processes aimed at solving technological problems
    \item can communicate and mediate research-based knowledge both orally and in writing
    \item is familiar with and can seek out leading international research within their specialist area.
    \item can work independently and reflect on own learning, academic development, and specialization
    \item masters technical problem-solving at a high level through cross-disciplinary teamwork, and can work with and manage all phases of a project – including preparation of timetables, design, solution, and documentation
\end{enumerate}

\end{document}